\subsection{Local Gauge Symmetry}
% This replaces the inadequate Definition 2.3 in the current paper

\begin{definition}[Gauge Transformations on Ledger States]
A gauge transformation $g: \mathbb{N} \to SU(3)$ acts on a configuration state $S$ by:
\begin{equation}
(g \cdot S)_n = U_g(n) S_n U_g^\dagger(n)
\end{equation}
where $U_g(n) \in SU(3)$ and the transformation preserves the ledger structure through:
\begin{equation}
d_n' = \text{Tr}(T^a U_g(n) d_n U_g^\dagger(n)), \quad c_n' = \text{Tr}(T^a U_g(n) c_n U_g^\dagger(n))
\end{equation}
with $T^a$ the Gell-Mann matrices.
\end{definition}

\begin{lemma}[Gauge Invariance of Cost Functional]
The cost functional $\mathcal{C}(S)$ is invariant under local gauge transformations:
\begin{equation}
\mathcal{C}(g \cdot S) = \mathcal{C}(S)
\end{equation}
for all $g \in \text{Map}(\mathbb{N}, SU(3))$.
\end{lemma}

\begin{proof}
The key observation is that the cost functional depends only on gauge-invariant combinations. For each level $n$:
\begin{align}
|d_n' - c_n'| + d_n' + c_n' &= |\text{Tr}(T^a U_g d_n U_g^\dagger) - \text{Tr}(T^a U_g c_n U_g^\dagger)| \\
&\quad + \text{Tr}(T^a U_g d_n U_g^\dagger) + \text{Tr}(T^a U_g c_n U_g^\dagger) \\
&= |d_n - c_n| + d_n + c_n
\end{align}
using the cyclic property of the trace and unitarity of $U_g$.
\end{proof}

\begin{definition}[Physical Gauge Sector]
The physical gauge sector $\mathcal{G}_{\text{phys}}$ consists of gauge-equivalence classes of states:
\begin{equation}
\mathcal{G}_{\text{phys}} = \{[S] : S \in \mathcal{G}\} / SU(3)_{\text{local}}
\end{equation}
where $[S]$ denotes the orbit of $S$ under gauge transformations.
\end{definition}

\begin{theorem}[Emergence of Color Charge]
The mod 3 structure of indices naturally selects states transforming non-trivially under the center $Z_3 \subset SU(3)$:
\begin{equation}
n \bmod 3 \neq 0 \iff [S_n] \text{ carries non-zero triality}
\end{equation}
This provides a gauge-invariant characterization of the physical sector.
\end{theorem}

\begin{proof}
The center $Z_3 = \{e^{2\pi i k/3} \mathbb{I} : k = 0,1,2\}$ acts on states by phases. States at level $n$ transform as:
\begin{equation}
Z_3: S_n \mapsto e^{2\pi i n/3} S_n
\end{equation}
Non-trivial transformation occurs precisely when $n \bmod 3 \neq 0$, establishing the connection between the discrete index structure and gauge quantum numbers.
\end{proof}

\begin{remark}[Gauss Law Constraints]
The physical Hilbert space $\mathcal{H}_{\text{phys}}$ is the kernel of the Gauss law operators:
\begin{equation}
G^a |S\rangle = 0, \quad a = 1,\ldots,8
\end{equation}
where $G^a$ generates infinitesimal gauge transformations. In the ledger formalism, this is automatically satisfied by working with gauge-invariant cost functionals.
\end{remark}

% Connection to standard formulation
\begin{proposition}[Relation to Wilson Lines]
The ledger entries $(d_n, c_n)$ can be interpreted as discretized Wilson line correlators:
\begin{equation}
d_n - c_n \sim \langle \text{Tr}(P e^{i\int A_\mu dx^\mu}) \rangle_n
\end{equation}
where the path integral is over loops at scale $n$. This provides the bridge to conventional gauge theory.
\end{proposition} 